\subsection{"以形补形"与壮阳食疗：传统说法的现代审视}

"吃啥补啥"的取象比类思维深植于民间食疗文化，围绕性功能的具体说法也最多。逐项过一遍证据：

\begin{itemize}
\item \textbf{生蚝与牡蛎}：这是少数"有一点真实基础"的说法——牡蛎含锌量确实突出，而锌与睾酮合成相关。但要补上后半句：锌的增益只对缺锌者成立，锌摄入正常的人额外补锌并无好处，长期超量反而干扰铜吸收、影响免疫；至于"当晚吃、当晚灵"，则完全是美好想象；
\item \textbf{韭菜}："壮阳草"的名声缺乏成分学依据——其锌含量在常见食物中并不突出，膳食纤维与硫化物对健康有益，但与性功能没有特异关联；
\item \textbf{各种"鞭"（动物阴茎与睾丸制品）}：主要成分是胶原蛋白与脂肪，不含特异性"壮阳物质"；野生动物制品还叠加法律与物种保护问题；
\item \textbf{鹿茸}：传统"峻补元阳"的代表，现代研究提示其含微量激素样物质与生长因子，临床证据薄弱；作为药材有采收合规问题，作为保健品则市场乱象多。
\end{itemize}

\textbf{"一滴精十滴血"与"伤肾"焦虑}：精液的主要成分是水、果糖、蛋白质与少量电解质，生成与排出由身体持续代谢完成，适度排精不损耗任何"元气"（参见本书遗精与自慰相关章节的讨论）。把正常生理消耗恐惧为"亏空"，由此产生的焦虑反而更容易引起性功能障碍——恐惧比"亏空"本身伤人。

\textbf{市场风险高于营养风险}：壮阳药酒与"秘方食疗"是非法添加的高发品类——监管抽检中检出西地那非类似物的案例屡见不鲜（相互作用与禁忌参见本章《壮阳中药与西药的相互作用》小节）。本书的立场与前文一致：传统食疗可以作为一种饮食文化与日常滋养来理解；涉及具体健康问题时，把决策权交给循证医学，把"以形补形"留在餐桌上，而不是药单里。
